\documentclass[12pt,a4paper]{article}
\usepackage[utf8]{vietnam}
\usepackage{amsmath,amssymb}
\usepackage{tabularx}
\usepackage{tikz}
\usepackage{xcolor}
\usepackage[top=0.6cm, bottom=0.6cm, left=0.6cm, right=0.6cm]{geometry}
\usepackage[dethi]{ex_test}

\setlength{\parskip}{0pt}

% Dòng kẻ chấm mờ đều đặn chuẩn viết tay (khoảng cách 7.2mm rộng rãi)
\newcommand{\dongke}[1]{%
  \begingroup%
  \setlength{\parskip}{0pt}%
  \baselineskip=20.5pt%
  \vspace{1.5mm}%
  \foreach \i in {1,...,#1}{%
    \noindent{\color{gray!60}\dotfill}\par%
  }%
  \vspace{1mm}%
  \endgroup%
}

% Định nghĩa ô bài tập tự luận trong bảng
\newcommand{\oBaiTap}[3][6]{%
  \vspace{1.5mm}%
  \noindent\textbf{#2} #3\par%
  \dongke{#1}%
}

% Câu trắc nghiệm 1 cột: 4 đáp án dàn hàng ngang trên 1 dòng
\newcommand{\cauTN}[6]{% \cauTN{số câu}{đề}{A}{B}{C}{D}
  \noindent\textbf{Câu #1.} #2\par\vspace{0.4mm}%
  \noindent%
  \makebox[0.25\linewidth][l]{\textbf{A.} #3}%
  \makebox[0.25\linewidth][l]{\textbf{B.} #4}%
  \makebox[0.25\linewidth][l]{\textbf{C.} #5}%
  \makebox[0.25\linewidth][l]{\textbf{D.} #6}\par\vspace{0.8mm}%
}

\begin{document}
\pagestyle{empty}

% =============================================================================
% ========================= TOÁN 6 - TUẦN 04 - BUỔI 01: MÃ ĐỀ 101 =============
% ========================= MẶT TRƯỚC (TRANG 1/2) ==============================
% =============================================================================

\noindent
\begin{tabularx}{\linewidth}{|X|c|}
\hline
\textbf{PHIẾU KIỂM TRA ĐỊNH KỲ TOÁN 6 - TUẦN 04 - BUỔI 01} & \textbf{Mã đề: 101 - Trang 1/2} \\
\textit{Chủ đề: Thứ tự thực hiện các phép tính} & \textit{Thời gian: 45 phút} \\
\hline
\end{tabularx}

\vspace{1mm}

\noindent
\begin{tabularx}{\linewidth}{|X|}
\hline
\rule[-3.5mm]{0pt}{10mm}\textbf{Họ và tên học sinh:} \hfill \textbf{Mã HS:} \makebox[3.5cm]{} \\
\hline
\end{tabularx}

\vspace{1mm}

\noindent
\begin{tabularx}{\linewidth}{|p{3.2cm}|X|}
\hline
\centering\textbf{ĐIỂM SỐ} & \centering\textbf{LỜI NHẬN XÉT CỦA THẦY CÔ} \tabularnewline
\hline
\rule{0pt}{1.3cm} & \tabularnewline
\hline
\end{tabularx}

\vspace{1.2mm}

\noindent\textbf{BẢNG ĐIỀN ĐÁP ÁN TRẮC NGHIỆM (Mỗi câu đúng 0,5 điểm):}
\vspace{0.8mm}

\noindent
\begingroup
\setlength{\tabcolsep}{3pt}
\renewcommand{\arraystretch}{1.35}
\begin{tabularx}{\linewidth}{|c|*{5}{>{\centering\arraybackslash}X|}}
\hline
\textbf{Câu} & \textbf{1} & \textbf{2} & \textbf{3} & \textbf{4} & \textbf{5} \\
\hline
\textbf{Đáp án} & & & & & \\
\hline
\end{tabularx}
\endgroup

\vspace{1.5mm}

\noindent\textbf{\large PHẦN I. TRẮC NGHIỆM KHÁCH QUAN (2,5 điểm)}
\vspace{1mm}

\noindent\textbf{Câu 1.} Thứ tự đúng khi thực hiện phép tính không có dấu ngoặc là:\par\vspace{0.4mm}%
\noindent%
\makebox[0.5\linewidth][l]{\textbf{A.} Nhân, chia $\to$ Lũy thừa $\to$ Cộng, trừ}%
\makebox[0.5\linewidth][l]{\textbf{B.} Lũy thừa $\to$ Nhân, chia $\to$ Cộng, trừ}\par\vspace{0.4mm}%
\noindent%
\makebox[0.5\linewidth][l]{\textbf{C.} Lũy thừa $\to$ Cộng, trừ $\to$ Nhân, chia}%
\makebox[0.5\linewidth][l]{\textbf{D.} Cộng, trừ $\to$ Nhân, chia $\to$ Lũy thừa}\par\vspace{0.8mm}%

\cauTN{2}{Thứ tự đúng khi thực hiện phép tính có các loại dấu ngoặc là:}{$\{ \} \to [ ] \to ( )$}{$[ ] \to ( ) \to \{ \}$}{$( ) \to \{ \} \to [ ]$}{$( ) \to [ ] \to \{ \}$}

\cauTN{3}{Kết quả của phép tính $20 - 2 \cdot 3 + 4$ là:}{$18$}{$58$}{$10$}{$14$}

\cauTN{4}{Giá trị của biểu thức $5 \cdot 2^2 - 18 : 3^2$ là:}{$16$}{$22$}{$18$}{$20$}

\cauTN{5}{Giá trị của biểu thức $100 : [20 + (15 - 10)]$ là:}{$5$}{$4$}{$10$}{$2$}

\vspace{1.2mm}

\noindent\textbf{\large PHẦN II. TỰ LUẬN (7,5 điểm)}

\noindent\textbf{Câu 1 (2,5 điểm). Thực hiện phép tính (hợp lý nếu có thể):}
\vspace{0.8mm}

\noindent
\begin{tabularx}{\linewidth}{|X|X|}
\hline
\oBaiTap[6]{a)}{$45 + 15 : 3 - 2^3$} & 
\oBaiTap[6]{b)}{$36 \cdot 45 + 36 \cdot 55 - 600$} \\
\hline
\oBaiTap[6]{c)}{$120 : [54 - (50 - 2^3 \cdot 2)]$} & 
\oBaiTap[6]{d)}{$160 : \{ 170 - [65 + (8 - 3)^2] \}$} \\
\hline
\end{tabularx}

\newpage
% =============================================================================
% ========================= TOÁN 6 - TUẦN 04 - BUỔI 01: MÃ ĐỀ 101 =============
% ========================= MẶT SAU (TRANG 2/2) ================================
% =============================================================================

\noindent
\begin{tabularx}{\linewidth}{|X|c|}
\hline
\textbf{PHẦN II. TỰ LUẬN (Tiếp theo) - TOÁN 6 TUẦN 04 BUỔI 01} & \textbf{Mã đề: 101 - Trang 2/2} \\
\hline
\end{tabularx}

\vspace{2.5mm}

\noindent\textbf{\large Câu 2 (2,5 điểm). Tìm số tự nhiên $x$, biết:}
\vspace{0.8mm}

\noindent
\begin{tabularx}{\linewidth}{|X|X|}
\hline
\oBaiTap[6]{a)}{$25 + 3(x - 2) = 40$} & 
\oBaiTap[6]{b)}{$120 - [45 - (x - 5)] = 85$} \\
\hline
\oBaiTap[6]{c)}{$3 \cdot 2^x - 5 = 19$} & 
\oBaiTap[6]{d)}{$2(x + 7) - 3^3 = 11$} \\
\hline
\end{tabularx}

\vspace{2.5mm}

\noindent\textbf{\large Câu 3 (2,5 điểm). Bài toán thực tế:}
\vspace{0.8mm}

\noindent
\begin{tabularx}{\linewidth}{|X|}
\hline
\vspace{0.5mm}
Bạn An đi nhà sách mua: $5$ quyển vở giá $12\,000$ đồng/quyển, $3$ cây bút bi giá $6\,000$ đồng/cây và $1$ hộp bút sáp màu giá $35\,000$ đồng. Nhà sách có chính sách giảm giá $15\,000$ đồng cho hóa đơn mua hàng có tổng giá trị từ $100\,000$ đồng trở lên.\par
\textbf{a) (1,25 điểm)} Hãy viết một biểu thức số thể hiện số tiền bạn An phải trả.\par
\textbf{b) (1,25 điểm)} Tính cụ thể số tiền bạn An phải trả cho nhà sách sau khi được giảm giá.\par
\vspace{0.4mm}
\noindent\textit{Bài giải:}\dongke{14}
\\ \hline
\end{tabularx}

\newpage
% =============================================================================
% ========================= TOÁN 6 - TUẦN 04 - BUỔI 01: MÃ ĐỀ 102 =============
% ========================= MẶT TRƯỚC (TRANG 1/2) ==============================
% =============================================================================

\noindent
\begin{tabularx}{\linewidth}{|X|c|}
\hline
\textbf{PHIẾU KIỂM TRA ĐỊNH KỲ TOÁN 6 - TUẦN 04 - BUỔI 01} & \textbf{Mã đề: 102 - Trang 1/2} \\
\textit{Chủ đề: Thứ tự thực hiện các phép tính} & \textit{Thời gian: 45 phút} \\
\hline
\end{tabularx}

\vspace{1mm}

\noindent
\begin{tabularx}{\linewidth}{|X|}
\hline
\rule[-3.5mm]{0pt}{10mm}\textbf{Họ và tên học sinh:} \hfill \textbf{Mã HS:} \makebox[3.5cm]{} \\
\hline
\end{tabularx}

\vspace{1mm}

\noindent
\begin{tabularx}{\linewidth}{|p{3.2cm}|X|}
\hline
\centering\textbf{ĐIỂM SỐ} & \centering\textbf{LỜI NHẬN XÉT CỦA THẦY CÔ} \tabularnewline
\hline
\rule{0pt}{1.3cm} & \tabularnewline
\hline
\end{tabularx}

\vspace{1.2mm}

\noindent\textbf{BẢNG ĐIỀN ĐÁP ÁN TRẮC NGHIỆM (Mỗi câu đúng 0,5 điểm):}
\vspace{0.8mm}

\noindent
\begingroup
\setlength{\tabcolsep}{3pt}
\renewcommand{\arraystretch}{1.35}
\begin{tabularx}{\linewidth}{|c|*{5}{>{\centering\arraybackslash}X|}}
\hline
\textbf{Câu} & \textbf{1} & \textbf{2} & \textbf{3} & \textbf{4} & \textbf{5} \\
\hline
\textbf{Đáp án} & & & & & \\
\hline
\end{tabularx}
\endgroup

\vspace{1.5mm}

\noindent\textbf{\large PHẦN I. TRẮC NGHIỆM KHÁCH QUAN (2,5 điểm)}
\vspace{1mm}

\noindent\textbf{Câu 1.} Khi biểu thức chỉ có các phép tính cộng và trừ (hoặc chỉ nhân và chia), ta thực hiện:\par\vspace{0.4mm}%
\noindent%
\makebox[0.5\linewidth][l]{\textbf{A.} Từ phải sang trái}%
\makebox[0.5\linewidth][l]{\textbf{B.} Thực hiện phép trừ trước}\par\vspace{0.4mm}%
\noindent%
\makebox[0.5\linewidth][l]{\textbf{C.} Từ trái sang phải}%
\makebox[0.5\linewidth][l]{\textbf{D.} Tùy ý chọn phép tính nào trước}\par\vspace{0.8mm}%

\cauTN{2}{Dấu ngoặc nào được ưu tiên thực hiện đầu tiên trong biểu thức có nhiều ngoặc?}{$( )$}{$[ ]$}{$\{ \}$}{Tùy ý dấu ngoặc}

\cauTN{3}{Kết quả của phép tính $30 - 3 \cdot 4 + 2$ là:}{$110$}{$16$}{$24$}{$20$}

\cauTN{4}{Giá trị của biểu thức $4 \cdot 3^2 - 32 : 2^3$ là:}{$36$}{$32$}{$30$}{$28$}

\cauTN{5}{Giá trị của biểu thức $120 : [30 + (20 - 10)]$ là:}{$4$}{$6$}{$3$}{$5$}

\vspace{1.2mm}

\noindent\textbf{\large PHẦN II. TỰ LUẬN (7,5 điểm)}

\noindent\textbf{Câu 1 (2,5 điểm). Thực hiện phép tính (hợp lý nếu có thể):}
\vspace{0.8mm}

\noindent
\begin{tabularx}{\linewidth}{|X|X|}
\hline
\oBaiTap[6]{a)}{$55 + 24 : 4 - 3^2$} & 
\oBaiTap[6]{b)}{$45 \cdot 37 + 45 \cdot 63 - 500$} \\
\hline
\oBaiTap[6]{c)}{$140 : [50 - (40 - 3^2 \cdot 2)]$} & 
\oBaiTap[6]{d)}{$240 : \{ 200 - [95 + (8 - 3)^2] \}$} \\
\hline
\end{tabularx}

\newpage
% =============================================================================
% ========================= TOÁN 6 - TUẦN 04 - BUỔI 01: MÃ ĐỀ 102 =============
% ========================= MẶT SAU (TRANG 2/2) ================================
% =============================================================================

\noindent
\begin{tabularx}{\linewidth}{|X|c|}
\hline
\textbf{PHẦN II. TỰ LUẬN (Tiếp theo) - TOÁN 6 TUẦN 04 BUỔI 01} & \textbf{Mã đề: 102 - Trang 2/2} \\
\hline
\end{tabularx}

\vspace{2.5mm}

\noindent\textbf{\large Câu 2 (2,5 điểm). Tìm số tự nhiên $x$, biết:}
\vspace{0.8mm}

\noindent
\begin{tabularx}{\linewidth}{|X|X|}
\hline
\oBaiTap[6]{a)}{$35 + 4(x - 3) = 55$} & 
\oBaiTap[6]{b)}{$130 - [55 - (x - 8)] = 90$} \\
\hline
\oBaiTap[6]{c)}{$2 \cdot 3^x - 5 = 49$} & 
\oBaiTap[6]{d)}{$3(x + 5) - 2^4 = 23$} \\
\hline
\end{tabularx}

\vspace{2.5mm}

\noindent\textbf{\large Câu 3 (2,5 điểm). Bài toán thực tế:}
\vspace{0.8mm}

\noindent
\begin{tabularx}{\linewidth}{|X|}
\hline
\vspace{0.5mm}
Bạn Bình đi nhà sách mua: $4$ cuốn truyện giá $18\,000$ đồng/cuốn, $2$ chiếc compa giá $15\,000$ đồng/chiếc và $1$ thước kẻ giá $8\,000$ đồng. Cửa hàng đang có voucher giảm giá $10\,000$ đồng cho hóa đơn mua hàng từ $100\,000$ đồng trở lên.\par
\textbf{a) (1,25 điểm)} Hãy viết một biểu thức số thể hiện số tiền bạn Bình phải trả.\par
\textbf{b) (1,25 điểm)} Tính cụ thể số tiền bạn Bình phải trả cho nhà sách sau khi dùng voucher.\par
\vspace{0.4mm}
\noindent\textit{Bài giải:}\dongke{14}
\\ \hline
\end{tabularx}

\newpage
% =============================================================================
% ========================= TOÁN 6 - TUẦN 04 - BUỔI 01: MÃ ĐỀ 103 =============
% ========================= MẶT TRƯỚC (TRANG 1/2) ==============================
% =============================================================================

\noindent
\begin{tabularx}{\linewidth}{|X|c|}
\hline
\textbf{PHIẾU KIỂM TRA ĐỊNH KỲ TOÁN 6 - TUẦN 04 - BUỔI 01} & \textbf{Mã đề: 103 - Trang 1/2} \\
\textit{Chủ đề: Thứ tự thực hiện các phép tính} & \textit{Thời gian: 45 phút} \\
\hline
\end{tabularx}

\vspace{1mm}

\noindent
\begin{tabularx}{\linewidth}{|X|}
\hline
\rule[-3.5mm]{0pt}{10mm}\textbf{Họ và tên học sinh:} \hfill \textbf{Mã HS:} \makebox[3.5cm]{} \\
\hline
\end{tabularx}

\vspace{1mm}

\noindent
\begin{tabularx}{\linewidth}{|p{3.2cm}|X|}
\hline
\centering\textbf{ĐIỂM SỐ} & \centering\textbf{LỜI NHẬN XÉT CỦA THẦY CÔ} \tabularnewline
\hline
\rule{0pt}{1.3cm} & \tabularnewline
\hline
\end{tabularx}

\vspace{1.2mm}

\noindent\textbf{BẢNG ĐIỀN ĐÁP ÁN TRẮC NGHIỆM (Mỗi câu đúng 0,5 điểm):}
\vspace{0.8mm}

\noindent
\begingroup
\setlength{\tabcolsep}{3pt}
\renewcommand{\arraystretch}{1.35}
\begin{tabularx}{\linewidth}{|c|*{5}{>{\centering\arraybackslash}X|}}
\hline
\textbf{Câu} & \textbf{1} & \textbf{2} & \textbf{3} & \textbf{4} & \textbf{5} \\
\hline
\textbf{Đáp án} & & & & & \\
\hline
\end{tabularx}
\endgroup

\vspace{1.5mm}

\noindent\textbf{\large PHẦN I. TRẮC NGHIỆM KHÁCH QUAN (2,5 điểm)}
\vspace{1mm}

\noindent\textbf{Câu 1.} Khi thực hiện phép tính trong biểu thức không có dấu ngoặc, thứ tự nào sau đây là đúng?\par\vspace{0.4mm}%
\noindent%
\makebox[0.5\linewidth][l]{\textbf{A.} Nhân, chia $\to$ Cộng, trừ $\to$ Lũy thừa}%
\makebox[0.5\linewidth][l]{\textbf{B.} Cộng, trừ $\to$ Nhân, chia $\to$ Lũy thừa}\par\vspace{0.4mm}%
\noindent%
\makebox[0.5\linewidth][l]{\textbf{C.} Lũy thừa $\to$ Cộng, trừ $\to$ Nhân, chia}%
\makebox[0.5\linewidth][l]{\textbf{D.} Lũy thừa $\to$ Nhân, chia $\to$ Cộng, trừ}\par\vspace{0.8mm}%

\cauTN{2}{Dấu ngoặc nào được thực hiện cuối cùng trong biểu thức có nhiều loại ngoặc?}{$\{ \}$}{$( )$}{$[ ]$}{Tùy ý thứ tự}

\cauTN{3}{Kết quả của phép tính $40 - 4 \cdot 6 + 5$ là:}{$221$}{$21$}{$15$}{$29$}

\cauTN{4}{Giá trị của biểu thức $6 \cdot 2^3 - 27 : 3^2$ là:}{$42$}{$41$}{$45$}{$51$}

\cauTN{5}{Giá trị của biểu thức $150 : [25 + (15 - 10)]$ là:}{$10$}{$6$}{$3$}{$5$}

\vspace{1.2mm}

\noindent\textbf{\large PHẦN II. TỰ LUẬN (7,5 điểm)}

\noindent\textbf{Câu 1 (2,5 điểm). Thực hiện phép tính (hợp lý nếu có thể):}
\vspace{0.8mm}

\noindent
\begin{tabularx}{\linewidth}{|X|X|}
\hline
\oBaiTap[6]{a)}{$65 + 18 : 2 - 2^4$} & 
\oBaiTap[6]{b)}{$54 \cdot 38 + 54 \cdot 62 - 400$} \\
\hline
\oBaiTap[6]{c)}{$160 : [58 - (42 - 2^3 \cdot 2)]$} & 
\oBaiTap[6]{d)}{$320 : \{ 180 - [75 + (11 - 6)^2] \}$} \\
\hline
\end{tabularx}

\newpage
% =============================================================================
% ========================= TOÁN 6 - TUẦN 04 - BUỔI 01: MÃ ĐỀ 103 =============
% ========================= MẶT SAU (TRANG 2/2) ================================
% =============================================================================

\noindent
\begin{tabularx}{\linewidth}{|X|c|}
\hline
\textbf{PHẦN II. TỰ LUẬN (Tiếp theo) - TOÁN 6 TUẦN 04 BUỔI 01} & \textbf{Mã đề: 103 - Trang 2/2} \\
\hline
\end{tabularx}

\vspace{2.5mm}

\noindent\textbf{\large Câu 2 (2,5 điểm). Tìm số tự nhiên $x$, biết:}
\vspace{0.8mm}

\noindent
\begin{tabularx}{\linewidth}{|X|X|}
\hline
\oBaiTap[6]{a)}{$45 + 3(x - 4) = 66$} & 
\oBaiTap[6]{b)}{$150 - [60 - (x - 6)] = 110$} \\
\hline
\oBaiTap[6]{c)}{$4 \cdot 3^x - 7 = 101$} & 
\oBaiTap[6]{d)}{$5(x + 4) - 2^5 = 38$} \\
\hline
\end{tabularx}

\vspace{2.5mm}

\noindent\textbf{\large Câu 3 (2,5 điểm). Bài toán thực tế:}
\vspace{0.8mm}

\noindent
\begin{tabularx}{\linewidth}{|X|}
\hline
\vspace{0.5mm}
Bạn Cúc đi nhà sách mua: $6$ quyển vở giá $10\,000$ đồng/quyển, $4$ cây bút gel giá $5\,000$ đồng/cây và $1$ bộ compa - êke giá $25\,000$ đồng. Nhà sách có chính sách giảm giá $15\,000$ đồng cho hóa đơn mua hàng có tổng giá trị từ $100\,000$ đồng trở lên.\par
\textbf{a) (1,25 điểm)} Hãy viết một biểu thức số thể hiện số tiền bạn Cúc phải trả.\par
\textbf{b) (1,25 điểm)} Tính cụ thể số tiền bạn Cúc phải trả cho nhà sách sau khi được giảm giá.\par
\vspace{0.4mm}
\noindent\textit{Bài giải:}\dongke{14}
\\ \hline
\end{tabularx}

\newpage
% =============================================================================
% ========================= TOÁN 6 - TUẦN 04 - BUỔI 01: MÃ ĐỀ 104 =============
% ========================= MẶT TRƯỚC (TRANG 1/2) ==============================
% =============================================================================

\noindent
\begin{tabularx}{\linewidth}{|X|c|}
\hline
\textbf{PHIẾU KIỂM TRA ĐỊNH KỲ TOÁN 6 - TUẦN 04 - BUỔI 01} & \textbf{Mã đề: 104 - Trang 1/2} \\
\textit{Chủ đề: Thứ tự thực hiện các phép tính} & \textit{Thời gian: 45 phút} \\
\hline
\end{tabularx}

\vspace{1mm}

\noindent
\begin{tabularx}{\linewidth}{|X|}
\hline
\rule[-3.5mm]{0pt}{10mm}\textbf{Họ và tên học sinh:} \hfill \textbf{Mã HS:} \makebox[3.5cm]{} \\
\hline
\end{tabularx}

\vspace{1mm}

\noindent
\begin{tabularx}{\linewidth}{|p{3.2cm}|X|}
\hline
\centering\textbf{ĐIỂM SỐ} & \centering\textbf{LỜI NHẬN XÉT CỦA THẦY CÔ} \tabularnewline
\hline
\rule{0pt}{1.3cm} & \tabularnewline
\hline
\end{tabularx}

\vspace{1.2mm}

\noindent\textbf{BẢNG ĐIỀN ĐÁP ÁN TRẮC NGHIỆM (Mỗi câu đúng 0,5 điểm):}
\vspace{0.8mm}

\noindent
\begingroup
\setlength{\tabcolsep}{3pt}
\renewcommand{\arraystretch}{1.35}
\begin{tabularx}{\linewidth}{|c|*{5}{>{\centering\arraybackslash}X|}}
\hline
\textbf{Câu} & \textbf{1} & \textbf{2} & \textbf{3} & \textbf{4} & \textbf{5} \\
\hline
\textbf{Đáp án} & & & & & \\
\hline
\end{tabularx}
\endgroup

\vspace{1.5mm}

\noindent\textbf{\large PHẦN I. TRẮC NGHIỆM KHÁCH QUAN (2,5 điểm)}
\vspace{1mm}

\noindent\textbf{Câu 1.} Khi thực hiện phép tính trong một biểu thức có các dấu ngoặc, thứ tự thực hiện đúng là:\par\vspace{0.4mm}%
\noindent%
\makebox[0.5\linewidth][l]{\textbf{A.} $( ) \to [ ] \to \{ \}$}%
\makebox[0.5\linewidth][l]{\textbf{B.} $[ ] \to ( ) \to \{ \}$}\par\vspace{0.4mm}%
\noindent%
\makebox[0.5\linewidth][l]{\textbf{C.} $\{ \} \to [ ] \to ( )$}%
\makebox[0.5\linewidth][l]{\textbf{D.} $( ) \to \{ \} \to [ ]$}\par\vspace{0.8mm}%

\cauTN{2}{Khi biểu thức chỉ có phép tính nhân và chia, ta thực hiện theo thứ tự:}{Từ phải sang trái}{Từ trái sang phải}{Nhân trước, chia sau}{Chia trước, nhân sau}

\cauTN{3}{Kết quả của phép tính $50 - 5 \cdot 4 + 8$ là:}{$188$}{$22$}{$38$}{$42$}

\cauTN{4}{Giá trị của biểu thức $3 \cdot 4^2 - 40 : 2^3$ là:}{$41$}{$43$}{$38$}{$45$}

\cauTN{5}{Giá trị của biểu thức $200 : [35 + (20 - 5)]$ là:}{$4$}{$5$}{$10$}{$8$}

\vspace{1.2mm}

\noindent\textbf{\large PHẦN II. TỰ LUẬN (7,5 điểm)}

\noindent\textbf{Câu 1 (2,5 điểm). Thực hiện phép tính (hợp lý nếu có thể):}
\vspace{0.8mm}

\noindent
\begin{tabularx}{\linewidth}{|X|X|}
\hline
\oBaiTap[6]{a)}{$75 + 32 : 8 - 2^3$} & 
\oBaiTap[6]{b)}{$28 \cdot 64 + 36 \cdot 28 - 800$} \\
\hline
\oBaiTap[6]{c)}{$210 : [48 - (50 - 2^4 \cdot 2)]$} & 
\oBaiTap[6]{d)}{$300 : \{ 220 - [95 + (7 - 2)^2] \}$} \\
\hline
\end{tabularx}

\newpage
% =============================================================================
% ========================= TOÁN 6 - TUẦN 04 - BUỔI 01: MÃ ĐỀ 104 =============
% ========================= MẶT SAU (TRANG 2/2) ================================
% =============================================================================

\noindent
\begin{tabularx}{\linewidth}{|X|c|}
\hline
\textbf{PHẦN II. TỰ LUẬN (Tiếp theo) - TOÁN 6 TUẦN 04 BUỔI 01} & \textbf{Mã đề: 104 - Trang 2/2} \\
\hline
\end{tabularx}

\vspace{2.5mm}

\noindent\textbf{\large Câu 2 (2,5 điểm). Tìm số tự nhiên $x$, biết:}
\vspace{0.8mm}

\noindent
\begin{tabularx}{\linewidth}{|X|X|}
\hline
\oBaiTap[6]{a)}{$50 + 4(x - 5) = 74$} & 
\oBaiTap[6]{b)}{$160 - [70 - (x - 10)] = 120$} \\
\hline
\oBaiTap[6]{c)}{$5 \cdot 2^x - 12 = 28$} & 
\oBaiTap[6]{d)}{$4(x + 6) - 5^2 = 27$} \\
\hline
\end{tabularx}

\vspace{2.5mm}

\noindent\textbf{\large Câu 3 (2,5 điểm). Bài toán thực tế:}
\vspace{0.8mm}

\noindent
\begin{tabularx}{\linewidth}{|X|}
\hline
\vspace{0.5mm}
Bạn Dũng đi nhà sách mua: $3$ cuốn sổ tay giá $25\,000$ đồng/cuốn, $5$ cây bút bi mực tím giá $6\,000$ đồng/cây và $1$ hộp gôm tẩy giá $15\,000$ đồng. Nhà sách có voucher giảm giá $20\,000$ đồng cho hóa đơn mua hàng có tổng giá trị từ $100\,000$ đồng trở lên.\par
\textbf{a) (1,25 điểm)} Hãy viết một biểu thức số thể hiện số tiền bạn Dũng phải trả.\par
\textbf{b) (1,25 điểm)} Tính cụ thể số tiền bạn Dũng phải trả cho nhà sách sau khi dùng voucher.\par
\vspace{0.4mm}
\noindent\textit{Bài giải:}\dongke{14}
\\ \hline
\end{tabularx}

\end{document}
